\documentclass[12pt]{article}
\usepackage[UTF8]{inputenc}
\usepackage{thaisupp}
\usepackage{enumerate}
\usepackage{etoolbox}
\AtBeginEnvironment{enumerate}{\itshape}
\begin{document}
\title{การเคลื่อนที่แบบโปรเจกไทล์ (เพิ่มเติม)}
\maketitle
\section*{คำถาม in class}
\begin{enumerate}
\def\labelenumi{\arabic*.}
\normalfont
\item \textbf{วัตถุถูกขว้างในแนวราบด้วยความเร็ว 8 m/s จากที่สูง 20 m วัตถุตกถึงพื้นห่างจากจุดปล่อยกี่เมตร (g = 10 m/s²)}
\begin{enumerate}[label=\alph*)]
\item 12 m
\item 16 m
\item 20 m
\item 24 m
\end{enumerate}
\item \textbf{วัตถุตกอิสระจากที่สูง 45 m กระทบพื้นด้วยความเร็วเท่าใด (g = 10 m/s²)}
\begin{enumerate}[label=\alph*)]
\item 20 m/s
\item 30 m/s
\item 40 m/s
\item 45 m/s
\end{enumerate}
\item \textbf{ในการเคลื่อนที่แบบโปรเจกไทล์ ณ ตำแหน่งใดของวิถีที่ความเร็วในแนวราบและแนวดิ่งมีขนาดเท่ากัน}
\begin{enumerate}[label=\alph*)]
\item จุดเริ่มต้น
\item จุดสูงสุด
\item จุดที่มุม 45° กับแนวราบ
\item จุดกระทบพื้น
\end{enumerate}
\item \textbf{ในการเคลื่อนที่แบบโปรเจกไทล์ ข้อใดต่อไปนี้ไม่จริง}
\begin{enumerate}[label=\alph*)]
\item ความเร่งคงที่ในแนวดิ่ง
\item ความเร็วคงที่ในแนวราบ
\item ความเร่งเปลี่ยนแปลงในแนวดิ่ง
\item วิถีเป็นรูปพาราโบลา
\end{enumerate}
\item \textbf{ขว้างวัตถุขึ้นในแนวดิ่งด้วยความเร็ว 25 m/s วัตถุจะขึ้นไปได้สูงสุดเท่ากับ (g = 10 m/s²)}
\begin{enumerate}[label=\alph*)]
\item 25 m
\item 31.25 m
\item 50 m
\item 62.5 m
\end{enumerate}
\item \textbf{ขว้างวัตถุขึ้นในแนวดิ่งด้วยความเร็ว 20 m/s วัตถุจะตกกลับลงมาถึงจุดเดิมใช้เวลากี่วินาที (g = 10 m/s²)}
\begin{enumerate}[label=\alph*)]
\item 2 s
\item 4 s
\item 6 s
\item 8 s
\end{enumerate}
\item \textbf{ยิงวัตถุในแนวมุม 37° ด้วยความเร็ว 50 m/s (sin 37° = 0.6, cos 37° = 0.8) ขนาดความเร็วที่จุดสูงสุดเท่ากับ}
\begin{enumerate}[label=\alph*)]
\item 30 m/s
\item 40 m/s
\item 50 m/s
\item 60 m/s
\end{enumerate}
\item \textbf{ยิงวัตถุในแนวมุม 60° กับแนวราบด้วยความเร็ว 20 m/s จากหน้าผาสูง 80 m วัตถุตกถึงพื้นที่ระยะทางเท่าใด (g = 10 m/s²)}
\begin{enumerate}[label=\alph*)]
\item 17.3 m
\item 34.6 m
\item 69.2 m
\item 138.4 m
\end{enumerate}
\item \textbf{ถ้าความเร็วเริ่มต้นเพิ่มขึ้น 2 เท่า ระยะไกลสุดในแนวราบจะเพิ่มขึ้นกี่เท่า (ยิงที่มุมเดิม)}
\begin{enumerate}[label=\alph*)]
\item 2 เท่า
\item 4 เท่า
\item 8 เท่า
\item 16 เท่า
\end{enumerate}
\item \textbf{ยิงวัตถุด้วยความเร็ว 40 m/s ที่มุม 53° (sin 53° = 0.8, cos 53° = 0.6) ความสูงสูงสุดเท่ากับ (g = 10 m/s²)}
\begin{enumerate}[label=\alph*)]
\item 32 m
\item 64 m
\item 80 m
\item 128 m
\end{enumerate}
\item \textbf{ยิงลูกบอลในแนวมุม 45° ด้วยความเร็ว 30 m/s วัตถุอยู่ในอากาศนานเท่าใด (g = 10 m/s²)}
\begin{enumerate}[label=\alph*)]
\item 2.12 s
\item 3.0 s
\item 4.24 s
\item 6.0 s
\end{enumerate}
\item \textbf{วัตถุถูกขว้างในแนวราบจากหน้าผาสูง 80 m ด้วยความเร็ว 5 m/s วัตถุจะตกกระทบพื้นที่เวลากี่วินาที (g = 10 m/s²)}
\begin{enumerate}[label=\alph*)]
\item 2 s
\item 4 s
\item 8 s
\item 16 s
\end{enumerate}
\item \textbf{นักกีฬาขว้างแหลนได้ระยะไกล 60 m ที่มุม 45° ถ้าต้องการให้ได้ระยะ 90 m โดยใช้มุมเดิม ต้องเพิ่มความเร็วเริ่มต้นกี่เท่า}
\begin{enumerate}[label=\alph*)]
\item 1.2 เท่า
\item 1.5 เท่า
\item √1.5 เท่า
\item 2 เท่า
\end{enumerate}
\item \textbf{ยิงวัตถุในแนวมุม 53° ด้วยความเร็ว 50 m/s (sin 53° = 0.8, cos 53° = 0.6) จงหาขนาดความเร็วเมื่อวัตถุอยู่ที่ความสูง 45 m (g = 10 m/s²)}
\begin{enumerate}[label=\alph*)]
\item 30.2 m/s
\item 35.6 m/s
\item 40.0 m/s
\item 50.0 m/s
\end{enumerate}
\item \textbf{ปล่อยวัตถุจากที่สูง 20 m พร้อมกับขว้างวัตถุอีกก้อนในแนวราบจากที่สูงเดียวกัน วัตถุทั้งสองจะตกพร้อมกันหรือไม่}
\begin{enumerate}[label=\alph*)]
\item ตกพร้อมกันเสมอ
\item วัตถุที่ถูกขว้างตกก่อน
\item วัตถุที่ถูกปล่อยตกก่อน
\item ขึ้นกับความเร็วของวัตถุที่ถูกขว้าง
\end{enumerate}
\item \textbf{ถ้าต้องการให้วัตถุไปถึงความสูงสูงสุด 45 m โดยใช้ความเร็วเริ่มต้นน้อยที่สุด ควรยิงที่มุมเท่าใด}
\begin{enumerate}[label=\alph*)]
\item 30°
\item 45°
\item 60°
\item 90°
\end{enumerate}
\item \textbf{ยิงวัตถุจากพื้นในแนวมุม 30° ด้วยความเร็ว 30 m/s ตกบนพื้นที่สูงกว่าจุดยิง 5 m ระยะไกลในแนวราบเท่ากับ (g = 10 m/s²)}
\begin{enumerate}[label=\alph*)]
\item 45 m
\item 67.5 m
\item 78.6 m
\item 90 m
\end{enumerate}
\item \textbf{ยิงวัตถุในแนวมุม 37° ด้วยความเร็ว 50 m/s จงหาขนาดความเร็วที่มุม 37° กับแนวราบ (g = 10 m/s²)}
\begin{enumerate}[label=\alph*)]
\item 25 m/s
\item 30 m/s
\item 40 m/s
\item 50 m/s
\end{enumerate}
\item \textbf{เครื่องบินบินขนานพื้นที่ความสูง 500 m ด้วยความเร็ว 540 km/hr ปล่อยลูกน้ำหนัก ลูกน้ำหนักจะตกห่างจากจุดปล่อยกี่เมตร (g = 10 m/s²)}
\begin{enumerate}[label=\alph*)]
\item 1,000 m
\item 1,500 m
\item 2,000 m
\item 2,500 m
\end{enumerate}
\item \textbf{ขว้างวัตถุจากหน้าผาสูง 45 m ในแนวมุม 30° กับแนวราบ วัตถุใช้เวลากี่วินาทีกว่าจะถึงพื้น (g = 10 m/s²)}
\begin{enumerate}[label=\alph*)]
\item 2.12 s
\item 3 s
\item 4.24 s
\item 6 s
\end{enumerate}
\item \textbf{ยิงวัตถุจากหน้าผาสูง 80 m ในแนวมุม 53° ด้วยความเร็ว 50 m/s (sin 53° = 0.8, cos 53° = 0.6) ความสูงสูงสุดจากพื้นดินเท่ากับ (g = 10 m/s²)}
\begin{enumerate}[label=\alph*)]
\item 80 m
\item 120 m
\item 160 m
\item 200 m
\end{enumerate}
\item \textbf{ยิงวัตถุจากพื้นในแนวมุม 60° ด้วยความเร็ว 20 m/s ตกลงพื้นที่สูงกว่าจุดยิง 5 m ระยะไกลในแนวราบเท่ากับ (g = 10 m/s²)}
\begin{enumerate}[label=\alph*)]
\item 17.3 m
\item 26 m
\item 30 m
\item 34.6 m
\end{enumerate}
\item \textbf{ยิงวัตถุด้วยความเร็วคงที่ที่มุม 15° และ 75° ข้อสรุปใดถูกต้อง}
\begin{enumerate}[label=\alph*)]
\item ระยะไกลเท่ากัน
\item มุม 15° ไกลกว่า
\item มุม 75° ไกลกว่า
\item ขึ้นกับความเร็วเริ่มต้น
\end{enumerate}
\item \textbf{ยิงวัตถุในแนวมุม 60° ด้วยความเร็ว 20 m/s กระทบพื้นที่มุม 30° กับแนวราบ ขนาดความเร็วกระทบพื้นเท่ากับ (g = 10 m/s²)}
\begin{enumerate}[label=\alph*)]
\item 10 m/s
\item 20 m/s
\item 30 m/s
\item 40 m/s
\end{enumerate}
\item \textbf{ระยะไกลสุดในแนวราบของการขว้างวัตถุที่มุม 30° กับแนวราบด้วยความเร็ว 40 m/s เท่ากับ (g = 10 m/s²)}
\begin{enumerate}[label=\alph*)]
\item 40 m
\item 80 m
\item 120 m
\item 138.6 m
\end{enumerate}
\item \textbf{วัตถุถูกขว้างขึ้นในแนวดิ่งด้วยความเร็ว 30 m/s จงหาความเร็วเมื่อวัตถุขึ้นไปได้ 20 m (g = 10 m/s²)}
\begin{enumerate}[label=\alph*)]
\item 10 m/s
\item 20 m/s
\item 30 m/s
\item 40 m/s
\end{enumerate}
\item \textbf{ยิงวัตถุในแนวมุม 37° กับแนวราบด้วยความเร็ว 50 m/s จงหาเวลาที่วัตถุขึ้นไปถึงจุดสูงสุด (g = 10 m/s², sin 37° = 0.6)}
\begin{enumerate}[label=\alph*)]
\item 2 s
\item 3 s
\item 4 s
\item 6 s
\end{enumerate}
\item \textbf{นักบาสเกตบอลยืนห่างจากห่วง 5 m ขว้างลูกบอลในแนวมุม 53° กับแนวราบด้วยความเร็ว 10 m/s ห่วงอยู่สูง 3.05 m ลูกบอลจะเข้าห่วงหรือไม่ (g = 10 m/s², sin 53° = 0.8, cos 53° = 0.6)}
\begin{enumerate}[label=\alph*)]
\item เข้าห่วงพอดี
\item ตกต่ำกว่าห่วง
\item ตกสูงกว่าห่วง
\item ข้อมูลไม่เพียงพอ
\end{enumerate}
\item \textbf{เครื่องบินบินขึ้นในแนวมุม 30° กับแนวราบด้วยความเร็ว 200 m/s ที่ความสูง 3,000 m ปล่อยลูกน้ำหนัก ลูกน้ำหนักจะตกห่างจุดปล่อยกี่เมตร (g = 10 m/s²)}
\begin{enumerate}[label=\alph*)]
\item 2,000 m
\item 4,000 m
\item 6,000 m
\item 8,000 m
\end{enumerate}
\item \textbf{ปล่อยวัตถุจากที่สูง h พร้อมกับขว้างวัตถุอีกก้อนขึ้นในแนวดิ่งด้วยความเร็ว v₀ ที่ระยะห่างจากจุดปล่อยต้องเป็นเท่าใดจึงจะตกพร้อมกัน}
\begin{enumerate}[label=\alph*)]
\item v₀² / g
\item v₀² / 2g
\item v₀² / 3g
\item v₀² / 4g
\end{enumerate}
\end{enumerate}
\end{document}